\section{An exact quotient operation for singleton dependency graphs}
\label{alg:setting}

The algebraic operation developed here removes many generators in one
simultaneous substitution. Its donors may involve arbitrarily many other
generators, with repeated letters and mixed signs. The compatibility
condition is an acyclic dependency graph. We first prove the quotient
statement for arbitrary presentations, and then bound the compressed
construction in terms of its actual input grammar. The rank is part of
the input throughout this section.

Let $X$ be the finite set of live generators, with $r=|X|$. Write $F(X)$
for the free group on $X$. A presentation has $m$ retained relator slots
$R_1,\ldots,R_m$, including any empty slots left by earlier operations.
The exact signed words in these slots are represented by roots of one
binary straight-line program $\mathcal G$. Its rules are signed letters,
binary concatenations, and a distinguished empty root. Let $M$ be the
number of nonempty grammar nodes reachable from the current roots.
The words are allowed to contain free cancellation. Their raw spellings,
rather than their reduced lengths, determine the occurrence conditions below.

\begin{definition}[A singleton DAG witness]
\label{alg:def-witness}
A witness chooses a set $C\subseteq X$ of $k$ distinct children and a
distinct donor slot $j(c)$ for each $c\in C$. The raw donor word must have
exactly one occurrence of either sign of $c$, giving a factorization
\[
 D_c:=R_{j(c)}=U_c c^{\epsilon_c}V_c,
 \qquad \epsilon_c\in\{1,-1\}.
\]
Define the right-hand side
\[
 W_c=(V_cU_c)^{-\epsilon_c}.
\]
For each selected child $d$ occurring in $W_c$, put an arrow $c\to d$.
The witness is valid when this graph on $C$ is acyclic. The surviving
generators are $A=X\setminus C$.
\end{definition}

The support of $W_c$ includes every absolute generator label occurring
in it. Signs and multiplicities affect the substituted word but not the
dependency graph. A donor consisting of $c^{\epsilon_c}$ alone is
permitted: its right-hand side is empty and it has no dependencies.
Dependencies on surviving generators also impose no restriction.

Choose a dependency order $c_1,\ldots,c_k$, so that an arrow
$c_i\to c_j$ implies $j<i$. Define words over $A$ recursively by
\[
 I_a=a\quad(a\in A),\qquad
 I_{c_i}=W_{c_i}\bigl[I_h/h:h\in X\bigr].
\]
Only already defined selected images are used in the second formula.
Negative letters receive inverse words, including reversal of their order.
These images define a homomorphism $\Phi:F(X)\to F(A)$.

\begin{theorem}[Simultaneous singleton quotient]
\label{alg:quotient}
For every valid singleton DAG witness,
\[
 \langle X\mid R_1,\ldots,R_m\rangle
 \cong
 \left\langle A\;\middle|\;
   \Phi(R_j),\ j\notin\{j(c):c\in C\}\right\rangle.
\]
Thus all retained roots can be mapped simultaneously, and exactly the
selected donor slots can then be emptied. The assertion holds for every
finitely presented group.
\end{theorem}

\begin{proof}
The equation $U_c c^{\epsilon_c}V_c=1$ is equivalent to
$c=(V_cU_c)^{-\epsilon_c}$; this also holds when either context is empty.
Let $N$ be the normal closure of the selected donors in $F(X)$. Every
donor is killed by $\Phi$, so $\Phi$ induces a homomorphism
$\overline\Phi:F(X)/N\to F(A)$. Inclusion of the surviving generators
induces a homomorphism $\iota:F(A)\to F(X)/N$.

We show in dependency order that $c_i=I_{c_i}$ in $F(X)/N$. The donor
equation gives $c_i=W_{c_i}$ there. Every selected letter in $W_{c_i}$
has an earlier index, so the induction hypothesis replaces it by its
image. This proves the assertion. The two homomorphisms are consequently
inverse on every surviving and selected generator. They are therefore
mutually inverse group homomorphisms. Finally quotient by the images
of all unselected relators. Under this isomorphism those images are
precisely the words $\Phi(R_j)$ stated in the theorem.
\end{proof}

The theorem justifies deletion of a donor even when its mapped raw word
is a nonempty spelling of the identity. No equality query is required for
that deletion. Every other original slot remains present; an unrelated
relator can impose a substantial condition after substitution.

\begin{remark}[Raw occurrences and source hypotheses]
\label{alg:raw-scope}
A signed exponent sum of $1$ or $-1$ does not establish the required
singleton condition: a word may contain a generator twice positively
and once negatively with noncommuting letters between the occurrences.
Conversely, an actual singleton donor needs no torsion-free hypothesis.
The ordinary, exponent-one case of a unit-coordinate power forest is
a special case of this construction. A proper-power donor with several
raw occurrences of its child is not automatically a singleton donor;
the existing root-evidence route remains a separate capability.
If a future extension infers $D_c=1$ from a checked proper-power donor
$D_c^q=1$, $|q|>1$, that inference will require an additional justification,
such as source-established torsion-freeness. It is separate from
Theorem~\ref{alg:quotient}.
\end{remark}

Acyclicity is also substantive. The two donors $ab$ and $ab^{-1}$ allow
the proposed choices $a$ and $b$ individually, but together create a
dependency cycle. Their presentation is $\Z/2\Z$: the first donor gives
$a=b^{-1}$ and the second gives $b^2=1$. Deleting both generators and
both donors would instead give the trivial group. An individual
singleton test therefore cannot replace the compatibility check.

\section{Linear grammar growth through shared contexts}
\label{alg:compression}

Compressed substitution is a well-established method in algorithmic
group theory; see Schleimer~\cite{Schleimer2008}. For fixed rank and a
fixed finite family of endomorphisms, Kapovich's quantitative layered
substitution statement constructs an SLP of size $O(M+t+1)$ for a
specified length-$t$ composition, with constants depending on that rank
and family~\cite[Proposition~3.2]{Kapovich2026}. Here the rank and the
donor words vary with the input. The following argument controls the
extraction and composition of these variable-sized definitions directly
from their shared source grammar.

\subsection{The disjointness of the selected occurrence paths}

For a selected child $c$, follow its unique occurrence from the donor
root through the binary grammar. At each concatenation, exactly one
child word contains $c$. This gives a path $H_c$ ending at a signed
terminal of $c$. The grammar can share that terminal and many other
nodes with other words; the path is defined by the occurrence in this
particular donor.

\begin{lemma}[Disjoint hole paths]
\label{alg:hole-paths}
For distinct selected children $c,d$, the paths $H_c$ and $H_d$ have
no grammar node in common. In particular,
\[
 \sum_{c\in C}|H_c|\le M.
\]
\end{lemma}

\begin{proof}
A node on $H_c$ expands to a word containing $c$. Since the donor
contains $c$ only once, every concatenation on the path has only one
child containing it, and a node containing $c$ cannot occur twice in
the donor's unfolded derivation. The path is thus well defined even
in a heavily shared grammar.

Suppose a node belonged to both $H_c$ and $H_d$. Its expansion would
contain both selected children. It occurs within $D_c$, so deleting
the unique occurrence of $c$ from that donor leaves an occurrence of
$d$. Hence $W_c$ contains $d$. The same reasoning within $D_d$ shows
that $W_d$ contains $c$. This creates the two-cycle $c\to d\to c$,
contradicting the witness condition. The paths are disjoint, and each
is a subset of the $M$ reachable source nodes, proving the sum bound.
\end{proof}

The lemma does not say that whole donor grammars are disjoint. Large
subwords over shared parents or survivors can occur in many donors.
It controls precisely the paths on which cutting out the chosen
occurrences requires new grammar structure. This distinction removes
the otherwise unnecessary factor of $k$ from a bound obtained by
charging a full grammar-height traversal independently to every donor.

\begin{lemma}[All contexts from one source grammar]
\label{alg:contexts}
The words $V_cU_c$, for all selected children, can be represented using
the source nodes and at most $M$ additional binary concatenation nodes.
Their construction has $O(M+k)$ grammar events and uses no reduction,
equality test, or expanded-word allocation.
\end{lemma}

\begin{proof}
Descend $H_c$ while accumulating its prefix and suffix. When the path
enters a right child, append the left sibling to the prefix. When it
enters a left child, prepend the right sibling to the suffix. Siblings
are reused by node reference. Their order is the order they have in
the original word, so the final contexts are exactly $U_c$ and $V_c$.
Each internal path node calls at most one context concatenation.
One final concatenation joins the suffix to the prefix.

There is one terminal node on each $H_c$. Thus the number of possible
new context nodes, including the final joins, is at most
\[
 \sum_{c\in C}(|H_c|-1)+k
 =\sum_{c\in C}|H_c|\le M.
\]
Empty contexts and interning can reduce this number.

The independent implementation obtains the two contexts by the exact
slices $[0,p)$ and $[p+1,|D_c|)$, where $p$ is the unique occurrence
position. Both boundary traversals follow $H_c$. At a path node, only
one of the two complementary slices can need a new concatenation:
which one depends on which child contains the hole. Fully included
sibling subwords are reused. The same bound therefore applies to this
alternative construction.
\end{proof}

\subsection{Compiling both orientations of every needed word}

Let $\mathcal H$ be the union of the source grammar and the context
nodes from Lemma~\ref{alg:contexts}. It has at most $2M$ nodes. For
each node $v$ use two virtual states $(v,+1)$ and $(v,-1)$, representing
its final substituted word and the inverse of that word. If $v=uv'$,
the positive state depends, in order, on $(u,+1),(v',+1)$; the negative
state depends on $(v',-1),(u,-1)$. An unselected terminal produces
itself or its inverse.

A terminal $c^\delta$ with $c\in C$ is redirected to the context
$V_cU_c$. In orientation $\sigma$, its context orientation is
$-\epsilon_c\delta\sigma$. This one sign formula handles both donor
signs, both occurrences of the generator in other words, and both
virtual orientations. It also avoids constructing and traversing a
separate inverse grammar for each already composed image.

\begin{lemma}[Acyclic signed compilation]
\label{alg:signed}
The augmented virtual graph is acyclic. Evaluating each needed virtual
state once produces the images $I_c^{\pm1}$ and all required mapped
relators with at most $4M$ reachable nonempty output grammar nodes.
\end{lemma}

\begin{proof}
Before terminal redirects, the graph follows source and context
concatenations and is acyclic. A directed cycle introduced by the
redirects must pass through selected terminals. Between two successive
selected-terminal redirects, a path starts in a defining context for
some child $c$ and reaches an occurrence of another selected child $d$.
Such a segment witnesses the dependency $c\to d$. A virtual cycle
would therefore give a directed closed walk, and hence a directed
cycle, in the dependency graph. The stipulated acyclicity excludes it.

Evaluate virtual states in a topological traversal. A terminal needs
at most one letter node, a concatenation at most one concatenation
node, and a redirect only an alias to its already evaluated state.
The defining equations and the sign convention show inductively that
these values are the desired substituted words. There are at most
$2|\mathcal H|\le4M$ virtual states. Each output node is the value of
one of these states, proving the node bound. States shared by several
images or relators are evaluated only once.
\end{proof}

All definitions are included among the traversal goals, even if no
unselected relator uses their children. This prevents a checker from
ignoring an invalid but unused part of a proposed schedule. Only after
the complete construction succeeds are the retained root array and
the live-generator set replaced. A failed resource check can leave
newly allocated arena nodes, while the represented presentation
remains the original one.

\begin{theorem}[Uniform size bound]
\label{alg:linear-size}
All node counts in this statement exclude the distinguished empty node.
A valid singleton DAG batch on a shared reachable source grammar of
$M$ nodes has a compressed output with at most $4M$ nonempty grammar
nodes and the original $m$ root slots. A persistent arena can perform
the construction by adding at most $5M$ nodes to its previous allocated
count. In particular, when no unreachable nodes were previously
allocated, its total allocated count is at most $6M$.
\end{theorem}

\begin{proof}
At most $M$ context nodes are allocated, followed by at most $4M$
nodes from signed compilation. The source nodes already existed.
Lemma~\ref{alg:signed} bounds the reachable output independently of
the source and temporary context nodes retained by the arena.
Previously allocated unreachable nodes are unchanged and must be
added to the total when applying an arena-wide memory cap.
\end{proof}

The constants are conservative; cached letters, empty contexts and
interning often reuse nodes. The substantive property is their
independence of dependency depth, rank, expanded donor lengths, and
the number of occurrences of a parent in another child's definition.

\section{An ordered certificate with small summaries}
\label{alg:ordered}

The version-eight move records an ordered list of pairs consisting of
an original donor slot and a chosen child. No substituted words,
occurrence positions, supports, or grammar node identifiers are trusted
as certificate data. The order itself is a topological witness.

Give every survivor rank zero and the $i$th chosen child rank $i$.
For every source node $v$, compute $q(v)$, the largest selected rank
in its expansion, and $t(v)$, its occurrence count capped at two.
The empty root has summary $(0,0)$ and a terminal of rank $h$ has
summary $(h,1)$. At a concatenation $v=uv'$, put
\[
 q(v)=\max\{q(u),q(v')\},
\]
\[
 t(v)=\min\!\left\{2,
 t(u)\mathbf 1_{q(u)=q(v)}+
 t(v')\mathbf 1_{q(v')=q(v)}\right\}.
\]
Counts at rank zero need not distinguish the surviving generators:
every selected donor must have a positive maximum rank.

\begin{proposition}[Verification by ranked maxima]
\label{alg:ranked-max}
An ordered list with distinct children and distinct donor slots is a
valid singleton DAG witness in the supplied order if and only if the
donor of its $i$th child has summary $(i,1)$.
Every valid singleton DAG admits an order satisfying these conditions.
\end{proposition}

\begin{proof}
The recurrence computes the maximum rank and the capped multiplicity
of that rank by induction over the source grammar. Only the $i$th
selected child has rank $i$. Summary $(i,1)$ says that this child occurs
exactly once and no later child occurs in its donor. Removing its
occurrence leaves only survivors and earlier children, so every
dependency is earlier. Conversely, those two occurrence conditions
immediately imply summary $(i,1)$. Finally every finite acyclic graph
has a dependency order, to which the equivalence applies.
\end{proof}

Each summary has constant cardinality and $O(\log(r+1))$ bits. The
check avoids storing an $r$-bit support mask at every grammar node.
The same maxima locate each unique occurrence: on its hole path,
exactly one child has the donor's maximum rank. The contexts can
therefore be independently sliced without an expanded occurrence scan.

The checker first verifies the record types, live generator membership,
slot ranges and distinctness conditions, and checks that source terminals
use live labels. It then checks each donor's ranked maximum against the
intact current roots before constructing that donor's contexts. No
presentation mutation is committed until every donor and the full
composition have passed. Boolean values are not accepted as integer
identifiers. All other relator slots are mapped and retained. These
checks connect the combinatorial witness to the
exact quotient theorem; a source-bound public verdict must additionally
reconstruct the diagram and replay the complete preceding proof prefix.

The producer uses a simpler order already present in the input: numeric
generator order. It caches, for each immutable node, its maximum
absolute label and that label's count capped at two. The first current
slot whose maximum occurs once supplies that maximum's donor. Chosen
children are then sorted increasingly. All their dependencies are
smaller, so the result is valid. Cached content summaries never cache
slot positions or a previous live-generator set; those are reconstructed
for each current presentation.

This numeric policy is incomplete. Let an anchor $a$ have a larger
numeric label than every leaf $c_i$, and take donors $a^q c_i$ with
$q\ge2$. Their maxima repeat, so the producer finds no pivot. Selecting
every $c_i$ is nevertheless a valid batch with a common surviving
anchor. A suitable order places the anchor first and exposes all
these donors. The general ordered checker accepts that witness.
Thus a restriction in discovery does not become a restriction on the
mathematical class of certificates it can verify.

\section{Encoding, bit operations, and phase composition}
\label{alg:complexity}

Put $N=M+r+m$, and let $\ell$ bound the bit length of generator labels.
After reindexing the reachable grammar, pointer fields require
$O(\log(N+2))$ bits. The output grammar, live-generator description and
retained root list consequently have canonical encoded length
\[
 O\!\left(N\bigl(\log(N+2)+\ell\bigr)\right).
\]
This bound describes the word representation itself. Cached arithmetic
metadata in an implementation can be larger.

In particular, a binary grammar with $Q$ nonempty nodes represents no
word longer than $2^Q$. Induction in topological order proves this:
a terminal has length one and each concatenation adds two lengths
of earlier nodes. The output word-length integers therefore have
$O(M+1)$ bits by Theorem~\ref{alg:linear-size}. The source, context and
compiled nodes involved in this batch can therefore require $O(M^2)$
bits of cached length metadata. Any metadata belonging to previously
allocated unreachable nodes is additional unchanged storage. Construction
operates on these integers; it does not pay for their represented
word lengths. For a retained root list $\mathcal R$, define
\[
 B(\mathcal R)=\max\!\left\{1,
   \left\lceil\log_2\!\left(1+
     \max\bigl(\{\len(R):R\in\mathcal R\}\cup\{0\}\bigr)
   \right)\right\rceil\right\}.
\]
Let $B_{\mathrm{in}}$ and $B_{\mathrm{out}}$ denote this quantity before
and after the batch, respectively, using the exact raw word lengths.
A separate elementary bound is
\[
 B_{\mathrm{out}}\le (k+1)B_{\mathrm{in}}+O(k+1),
\]
To see this, put $L=\max(1,\max_j\len(R_j))$ for the original roots.
Every defining word has at most $L$ letters. Defining the images in
dependency order multiplies the largest image length by at most $L$
at each step, and mapping an original root introduces one further
factor of at most $L$. The final root length is at most $L^{k+1}$.
The direct grammar bound is frequently sharper and suffices for
polynomial bit complexity.

Ranked summary evaluation, context extraction and signed compilation
have $O(N)$ grammar or slot events. Sorting the producer's $k$ selected
labels adds $O(k\log(k+1))$ comparisons; the maintained work counter
charges this sorting explicitly. Thus linear grammar growth does not
mean exactly linear producer work. Graph indexing and the source
postorder also have their ordinary implementation costs. A deterministic
implementation can use comparison-based indexing, and a Python
implementation uses its usual dictionary behavior.

The verifier's rank summaries have $O(\log(r+1))$ bits; the producer's
numeric-label maxima have $O(\ell)$ bits. Generator comparisons have
$O(\ell)$-bit operands, and position or length additions have
$O(M+1)$-bit operands. The number of these operations is polynomial
in $N$. This gives a deterministic polynomial bit bound in the encoded
input parameters, with an absolute degree independent of the rank
and DAG depth. Under the straightforward arithmetic representation,
the possible quadratic total cost of length additions is already a
conservative polynomial allowance. The abstract work counter is a
cooperative resource measure, rather than a claim to count those bit
operations exactly.

An internal application helper also accepts unsorted acyclic schedules
through a general dependency-mask fallback. Its polynomial mask cost
is separate from the public ordered-certificate route. Similarly,
literal replay pays for actual expanded output and can exhaust its
allocation allowance on a batch whose compressed construction is small.

\begin{theorem}[Composition of raw singleton batches]
\label{alg:phases}
Suppose $d$ successive operations consist solely of valid singleton DAG
contractions. If the initial encoded parameters, including any retained
arena storage, are polynomially bounded by $N_0$, the total construction
and ordered-verification cost
has a bound
\[
 \operatorname{poly}(N_0)\,2^{O(d)}.
\]
The constants do not depend on the ranks or dependency depths of
the individual batches.
\end{theorem}

\begin{proof}
Each contraction multiplies the reachable grammar bound by at most an
absolute constant, and a persistent arena obeys the conservative
recurrence $A_{i+1}\le6A_i$ for its nonempty allocated node count. Rank never
increases and the number of retained slots is fixed. Hence all grammar
and metadata parameters after $i$ phases are bounded by a polynomial
in $N_0$ times $6^{O(i)}$. Substitute this into the fixed polynomial
per-phase bound and sum over $i\le d$. The resulting geometric bound
has the stated form.
\end{proof}

From polynomial-size diagram data, $d=O(\log n)$ such phases would
have polynomial total raw-contraction cost; a polylogarithmic bound
on $d$ would give a quasipolynomial bound. The theorem does not bound
the discovery of useful new donors, the normalization work between
contractions, or the length of a search using other operations.
It is a local composition statement. For example, a later raw
singleton need not occur in its original slot: from slots $xg$ and
$x$, eliminating $x$ using the first exposes $g$ in the second.
Simply collecting the original slot--child pairs does not fuse that
history into one initial witness.

\section{The difficulty of choosing the largest batch}
\label{alg:selection}

Efficient verification of a supplied schedule leaves a separate
optimization problem. The following elementary reduction identifies
a concrete limit on exact unrestricted batch selection. Its source
is the classical independent-set problem, equivalently the complement
of vertex cover in the standard NP-completeness framework
of Karp~\cite{Karp1972}.

\begin{definition}[Singleton DAG selection]
\label{alg:selection-problem}
Given a finite list of explicit relator words and an integer $b$, ask
whether there is a valid singleton DAG witness selecting at least
$b$ distinct children. Both the donor assignment and a valid order
may be chosen. The compressed variant supplies the words by an SLP.
\end{definition}

\begin{theorem}[Selection is NP-complete]
\label{alg:selection-hardness}
Singleton DAG selection is NP-complete. NP-hardness already holds for
explicit positive words in which each relator has exactly one
singleton generator. Thus it also holds when the candidate child of
each donor is fixed in advance.
\end{theorem}

\begin{proof}
A certificate consists of the chosen distinct slots and children in
dependency order. Its size is polynomial in the input. Direct
occurrence and order checks verify it for explicit words; the ranked
summary verifier does so for compressed words. Both decision problems
therefore belong to NP.

For hardness, let $H=(V,E)$ be a finite simple undirected graph. Use
one generator $x_v$ for each vertex and one relator
\[
 D_v=x_v\prod_{u\in N_H(v)}x_u^2
 \qquad(v\in V),
\]
where each neighborhood is written in any fixed order. The total
explicit length is $|V|+4|E|$. The generator $x_v$ occurs once in
$D_v$, whereas every other generator present occurs twice. Thus
$x_v$ is the only possible singleton child for that donor.

A chosen set of donors corresponds to a vertex set $S\subseteq V$.
For adjacent selected vertices $u,v$, donor $D_u$ contains $x_v$ and
donor $D_v$ contains $x_u$. Their dependency graph contains the
two-cycle $x_u\to x_v\to x_u$. Therefore every valid selected set
is independent in $H$. Conversely, if $S$ is independent, none of
its selected donors contains another selected child. All its other
letters are survivors, so its dependency graph has no edges and is
valid. The largest possible batch has size exactly $\alpha(H)$.
Using the same threshold $b$ gives a polynomial reduction from
independent set and proves the claim.
\end{proof}

The reduction concerns arbitrary current-word presentations. It does
not establish hardness for Wirtinger presentations, for states known
to arise from a validated knot, or for the numeric-order planner.
Nor does it imply that finding a useful smaller batch is difficult
on a particular geometric family. It explains why maximum-cardinality
selection should be a distinct research objective rather than an
unproved property of a greedy scan.

\subsection{The recognition scope of a terminal batch}

If a checked batch leaves one generator $a$ and every retained mapped
relator has exponent sum zero at $a$, the resulting presentation is
$\langle a\mid\varnothing\rangle\cong\Z$. At rank one, exponent
sum completely determines a group word. At larger rank the same
test is insufficient: a commutator has zero exponent vector and can
still impose a nontrivial relation.

For a source-established knot group, the standard cyclic-group
characterization of the unknot supplies the final topological
conclusion. The construction theorem makes this a polynomial
positive route whenever the required source preparation and prefix
replay also have polynomial bounds. It supplies no claim that every
unknot reaches such a state efficiently.

The final public search policy invokes the operation only when it
finds at least two pivots and exactly $r-1$ selected children. Otherwise
the current presentation and the legacy operation order remain
unchanged. This terminal guard separates the valid general algebraic
operator from a policy for applying partial batches during heuristic
search. An algebraic quotient preserves the group while potentially
changing later normalization costs and search choices.

Relative to the same compressed search with its other options held
fixed, the guard gives logical preservation of positive successes
when resources are uncapped: it either reaches a checked rank-one
terminal or continues through the old presentation states. This
comparison does not identify the compressed search with the public
API's default explicit backend; enabling the option also selects
compressed search. Planning overhead still consumes real time and the shared
work allowance, so the assertion is not fixed-budget performance
dominance. The general all-DAG quotient and verifier remain useful
for future policies whose exposure and downstream progress can be
analyzed separately.
